% LTeX: language=de

\subsubsection{Aufgaben}

\begin{Exercise}{S. 195/16}{}
\begin{enumerate}[leftmargin=*]
    \item [16.]{
        13 Spiele \'{a} 3 Möglichkeiten\\[1.0ex]
        $3^{13}=1594323$ mögliche Ergebnisse
    }
\end{enumerate}
\end{Exercise}
\pagebreak

\begin{Exercise}{S. 195/18}{}
\begin{enumerate}[leftmargin=*]
    \item [18.]{
        Insgesamt stehen 6 verschiedene Kugeln zur Verfügung: B, G, R, S, V, W\\[1.0ex]
        \begin{tabularx}{0.5\linewidth}{@{}l | X | X | X @{}}
            & a) & b) & c) \\[1.0ex]
            1 & $5!$ & $4!$ & $3!$ \\[1.5ex]
            2 & $\dfrac{5!}{2!}$ & $\dfrac{4!}{2!}$ & $\dfrac{3!}{2!}$\\[1.5ex]
            3 & $6^{3}$ & $6^{2}$ & $6^{1}$\\
        \end{tabularx}
    }
\end{enumerate}
\end{Exercise}

\begin{Exercise}{S. 196/4}{}
\begin{enumerate}[leftmargin=*]
    \item [4.]{
        20 Kunden; 2 Möglichkeiten (schmeckt / schmeckt nicht)
        \begin{enumerate}[label=\alph*),leftmargin=*]
            \item[a)] {
                Kombination ohne Reihenfolge: \quad \begin{minipage}[t]{\linewidth}
                    $n = 20; \quad k = 15$\\[1.5ex]
                    $N = \dfrac{20!}{(20-15)! \cdot 15!} = 15504$\\
                \end{minipage}
                }
            \item[b)] {
                3 Kunden, denen es nicht schmeckt, bei 20 Kunden hintereinander anzuordnen erlaubt $20 - 3 + 1 = 18$ Möglichkeiten. Spätestens beim 18. Kunden muss die Serie beginnen.\\[1.5ex]
                $N = 18$\\
                }
            \item[c)] {
                Kombination ohne Reihenfolge: \quad \begin{minipage}[t]{\linewidth}
                    $n = 20; \quad k_{1} = 0; \quad k_{2} = 1; \quad k_{3} = 2$\\[1.5ex]
                    $N = \displaystyle \binom{20}{0} + \binom{20}{1} + \binom{20}{2} = 211$\\
                \end{minipage}
                }
        \end{enumerate}
    }
\end{enumerate}
\end{Exercise}

\begin{Exercise}{S. 196/7}{}
\begin{enumerate}[leftmargin=*]
    \item [7.]{
        24 Männer; 20 Frauen\\[1.0ex]
        Auswahl von 3 Männer und 2 Frauen
        \begin{itemize}[leftmargin=*]
            \item Männer: \quad \begin{minipage}[t]{\linewidth}
                Kombination ohne Reihenfolge; \quad $n = 24; \quad k = 3$\\[1.5ex]
                $N_{M} = \displaystyle \binom{24}{3} = 2024$\\
            \end{minipage}
            \item Frauen: \quad \begin{minipage}[t]{\linewidth}
                Kombination ohne Reihenfolge; \quad $n = 20; \quad k = 2$\\[1.5ex]
                $N_{F} = \displaystyle \binom{20}{2} = 190$\\
            \end{minipage}
        \end{itemize}
        Alle Möglichkeiten, die bei der Wahl entstehen können: $N = N_{M} \cdot N_{F} = 2024 \cdot 190 = 384560$\\
    }
\end{enumerate}
\end{Exercise}
\pagebreak

\begin{Exercise}{S. 196/8}{}
\begin{enumerate}[leftmargin=*]
    \item [8.]{
        Urnenmodell: \quad \begin{minipage}[t]{0.75\linewidth}
            \begin{itemize}[leftmargin=*]
                \item ohne zurücklegen (nochmal das gleiche Müsli beläuft sich auf eine schon berücksichtigte Kombination)
                \item ohne Reihenfolge (Im Müsli wird eh alles vermischt)
            \end{itemize}
            $n = 20; \quad k_{min} = 0; \quad k_{max} = 20$\\[1.5ex]
            $N = \displaystyle \sum_{k=0}^{20} \binom{20}{k} = 2^{20} = 1048576$\\
        \end{minipage}\\[1.0ex]
        Alternativ: \quad \begin{minipage}[t]{0.75\linewidth}
            20 Entscheidungen mit jeweils 2 Möglichkeiten (ja / nein)
            \begin{itemize}[leftmargin=*]
                \item mit zurücklegen (nächste Sorte Müsli wieder die gleichen Möglichkeiten: nehmen oder nicht nehmen)
                \item mit Reihenfolge (Unterscheidung der verschiedenen Sorten)
            \end{itemize}
            $n = 2$; \quad $k = 20$\\[1.5ex]
            $N = 2^{20} = 1048576$\\
        \end{minipage}
    }
\end{enumerate}
\end{Exercise}